\documentclass[10pt,a4paper]{article}

\usepackage[margin=1.05in]{geometry}
\usepackage{setspace}
\usepackage{amsmath,amssymb,amsthm}
\usepackage{graphicx}
\usepackage{booktabs}
\usepackage{longtable}
\usepackage{array}
\usepackage{caption}
\usepackage{enumitem}
\usepackage{titlesec}
\usepackage{microtype}
\usepackage{hyperref}
\usepackage{fancyhdr}

\hypersetup{colorlinks=true,linkcolor=black,urlcolor=black,citecolor=black}
\captionsetup{font=small,labelfont=bf,skip=4pt,justification=raggedright}

\titleformat{\section}{\normalsize\bfseries}{}{0em}{}
\titleformat{\subsection}{\small\bfseries}{}{0em}{}
\titleformat{\subsubsection}{\small\itshape}{}{0em}{}
\titlespacing*{\section}{0pt}{1.35em}{0.5em}
\titlespacing*{\subsection}{0pt}{0.95em}{0.35em}

\theoremstyle{plain}
\newtheorem{proposition}{Proposition}
\newtheorem{lemma}{Lemma}
\newtheorem{theorem}{Theorem}
\newtheorem{corollary}{Corollary}
\theoremstyle{definition}
\newtheorem{definition}{Definition}
\theoremstyle{remark}
\newtheorem{observation}{Observation}
\newtheorem{remark}{Remark}

\setstretch{1.14}
\setlist[itemize]{leftmargin=1.15em,itemsep=0.1em,topsep=0.2em}
\setlist[enumerate]{leftmargin=1.15em,itemsep=0.1em,topsep=0.2em}

\graphicspath{{drawings/}}

\pagestyle{fancy}
\fancyhf{}
\fancyhead[L]{\small\itshape Elements of Position}
\fancyhead[R]{\small\itshape I.\ The Closed Unit}
\fancyfoot[C]{\small\thepage}
\renewcommand{\headrulewidth}{0pt}

\begin{document}

\begin{center}
\vspace*{2.2em}
{\large Elements of Position}\\[1.6em]
{\LARGE\bfseries I.\ The Closed Unit}\\[2.2em]
{\normalsize Justin Erholtz}\\[0.35em]
{\small 12 September 2026}
\end{center}

\vspace{2.4em}

\begin{center}
{\large Where is 1?}
\end{center}

\vspace{1.8em}

\noindent
Where is \(1\)? Point to it.

You cannot, and not because the mark is missing.
One is not an object on the table. It is a concept.
A finished hold --- the interior of which can be cut without end.
This work lives in that double fact: the unit closes, and the cut does not.

If a whole can be split, two old means already know when the split is even.
They meet only when the parts are equal, and on the unit that meeting is the midpoint --- the only cut that is the same amount from both ends.
Inversion through the same unit sends every pair to another pair whose product is still one.
A circle of diameter one is asked only later, and only as a question of closure: when does the figure finish a turn?
The answer is an integer.
The walk that made the turn has a length, and that length stands to its inverse as parts of the same whole.

Nothing after that is a new kind of object.
The count repeats the unit.
The drawings are one building seen from several stations.
Infinity is already inside the first interval; it is not a later chapter and not a last tick.

\vspace{1.4em}

\noindent
\textit{A note.}
The argument stays inside the questions that produced it.
Claims that can be proved here are proved by school geometry.
The drawings are one full architectural style set, built in full by \texttt{Elements of Position - Drawings.py} from the same mathematical identities as the text.

\vspace{0.9em}
\noindent
\textit{A note on names used here.}
Closed unit --- \(1\), read as a finished whole.
Primary ratio --- \(1/2\), the one self-equal cut.
Prime ratio --- \(2\), the inverse of that cut (a name, not a derivation of primality).
Own-origin --- the height \(n/2\) at finished count \(n\).
Street --- an independent rigid process (the circle's walk, inversion, the midpoint cut) that owns its own ratios.
Unit --- a finished comparison between streets.
Vacuum / trace --- a relation named in one stroke, versus the same relation walked with cost.
Master sequence --- the successor count \(0,1,2,3,\ldots\), which cannot be outrun.

\newpage

\tableofcontents
\newpage

% =====================================================================
\section{One}
% =====================================================================

Where is \(1\)?
Point to it.

It is an idea, and it is real in the abstract sense.
There is no object sitting on the table of which one can say \emph{that is one} and stop.
What there is, is one \emph{of} a thing: one stone, one step, one finished hold.
In its most basic sense, then, one is not a specimen.
It is a closed object --- a unit that has completed itself.

That already splits.
Between \(0\) and \(1\) the cut can be repeated without end.
The interior of the unit is already endless; in that reading one never closes.
In the other reading one \emph{is} the close: the hold is finished, and that finished hold is what the word names.
Both readings are true of the same mark.
The work begins at that tension and does not dissolve it.

The endless interior is not decoration, and it is not a later chapter.
It is the location of infinity, named now so it cannot return as a mystical last index.
Infinity is already inside \((0,1)\).
It is THERE.
Walk there and you are HERE.
The same place, other writing.

If one is a finished whole, it can be split, and the parts can be compared.
The arithmetic mean and the geometric mean of two nonnegative pieces already know when a split is even: they meet only if the pieces are equal.
On a whole of \(2\) that meeting is \(1\).
On a whole of \(1\) it is \(\tfrac12\).
That \(\tfrac12\) is the only cut that is the same amount from both ends.
Inversion through the unit sends the pair to another pair whose product is still \(1\).
The lock of that inversion is the same \(1\) we have not yet pointed to.

None of this mentions a circle yet.

The circle arrives only when one asks a different question, still elementary: if the diameter is exactly \(1\), when does the figure close?
The answer is an integer --- the same \(1\).
The walk that made the close has length \(\pi\).
That length may be written with its inverse as a ratio of parts of the unit: \(\pi\) and \(1/\pi\), whose product and geometric mean are \(1\).
Equal, here, does not mean the decimals match.
It means the parts name each other through the whole.

Two faces of that close stay in force from here on.
The frozen interior of a finished circle: the unit holding its parts.
The live walk of the same circle: the running count at rate \(1\).
At \(n=2\) the live mean of a diameter-\(2\) walk is \(\pi\).
Both readings are true.
Neither is allowed to erase the other.

The count that follows is only that unit asked again, integer by integer, with no leap to a last mark.
Each \(n\) is another even split of the same whole, and also, by inversion, another size.
AM and GM still share a point.
The point is still a ratio of parts.
The complication is already here, and it does not get more exotic later: a close that is an integer, and a cut that does not end.

% ---------------------------------------------------------------------
\section{What the literature already holds}
% ---------------------------------------------------------------------

The pieces are old.
That is worth saying before any assembly looks like a discovery of new constants.

AM--GM for two terms is the geometric-mean theorem wearing an inequality.
A semicircle on a segment \(a+b\), radius the arithmetic mean, altitude to the division point the geometric mean, is the standard picture.
Thales makes the large triangle right-angled.
Similar right triangles give \(\text{altitude}^{2}=ab\).
The altitude cannot exceed the radius.
Algebraically the same fact is \((\sqrt{a}-\sqrt{b})^{2}\ge 0\).
None of that is new, and none of it is used below as ornament.

A circle measured as circumference over diameter, or as diameter over circumference, is equally old.
Inversion in a circle --- \(r\mapsto k^{2}/r\), and on the unit circle \(r\mapsto 1/r\) --- is a classical transformation: circles to circles, the unit circle fixed pointwise.
The pair \(\bigl(n,1/n\bigr)\) is that transformation on the ray.
The pair \(\bigl(\pi,1/\pi\bigr)\) is the same transformation applied to the walk of the paper circle.

Archimedes, in \emph{Measurement of a Circle}, compares a circle to inscribed and circumscribed regular polygons and bounds the ratio of circumference to diameter between \(3\tfrac{10}{71}\) and \(3\tfrac{1}{7}\) by carrying the comparison to the \(96\)-gon.
The spoke triangles used later here are those triangles.
The sandwich \(n\sin(\pi/n)<\pi<n\tan(\pi/n)\) is that comparison at a single finished \(n\).
Archimedes continues the count in order to squeeze a number.
The present work stops the count in order to keep the close visible as an integer and the inverse parts visible as a pair whose product is \(1\).

What is missing as one elementary statement is the joint itself: the close of the diameter-\(1\) circle named as the integer \(1\), and that same close written as inverse parts \(\pi\) and \(1/\pi\) whose product and geometric mean are that \(1\).
The later integers, the half-walk, the two motions of one structure --- those are the same joint asked again, and they can wait until they are built.

The reason the joint is under-said is ordinary.
Once \(\pi\) is named as circumference over diameter, the inverse reading looks like a footnote, and the geometric mean of any inversion pair is always \(1\), so specialising to this circle looks too cheap to title.
The work ran toward computing \(\pi\) and toward identifying the area-constant with the circumference-constant, and each half of that work was already obvious, which is how the hinge was left unnamed.

From here the writing stays inside that hinge.
It does not claim a new constant.
It names a joint the constants already share.

% ---------------------------------------------------------------------
\section{Two parts of a finished whole}
% ---------------------------------------------------------------------

The question was never just where \(1\) sits --- it was whether a finished thing can still be cut.
The two means answer that for any pair before a circle is even in the room.
If a finished whole can be split, the first honest question is when the split is even, and the arithmetic mean and the geometric mean of two nonnegative pieces already know the answer.

\subsection{The inequality}

\begin{proposition}
\label{prop:amgm}
For nonnegative \(a,b\),
\[
\frac{a+b}{2}\ge\sqrt{ab},
\]
with equality if and only if \(a=b\).
\end{proposition}

\begin{proof}
\(\bigl(\sqrt{a}-\sqrt{b}\bigr)^{2}\ge0\) expands to \(a+b\ge 2\sqrt{ab}\).
Divide by \(2\).
Equality holds precisely when \(\sqrt{a}=\sqrt{b}\).
\end{proof}

\subsection{A fixed sum}

Fix the sum at \(2\) and write the pair as \(x\) and \(2-x\), with \(x\in[0,2]\).
The arithmetic mean is constantly \(1\).
The geometric mean is \(\sqrt{x(2-x)}\), which is the upper semicircle of radius \(1\) centred at \((1,0)\).
The line \(y=1\) sits on or above that semicircle and touches it only at the peak, which is the one place the two parts are equal.

\begin{proposition}
On the pair \((x,2-x)\) one has \(\mathrm{AM}=1\), and \(\mathrm{AM}=\mathrm{GM}\) if and only if \(x=1\).
\end{proposition}

\begin{proof}
\(\mathrm{AM}=(x+(2-x))/2=1\).
By Proposition~\ref{prop:amgm}, equality holds if and only if \(x=2-x\), hence \(x=1\).
\end{proof}

The meeting value is \(1\), not \(\tfrac12\).
It is a half only in this sense: the whole being split is \(2\), and the parts are equal.

On the unit interval the same fact scales without changing its kind.
For the pair \((x,1-x)\) one has \(\mathrm{AM}=\tfrac12\), and equality only at \(x=\tfrac12\).
The common point is the midpoint of the unit, and the ratio of parts is again \(1:1\).
That is the only cut on a finished segment that is the same amount from both ends.

\begin{figure}[ht]
\centering
\includegraphics[width=0.48\textwidth]{drawing_01_amgm_sum2.png}\hfill
\includegraphics[width=0.48\textwidth]{drawing_02_amgm_sum1.png}
\caption{AM--GM on a finished whole of \(2\), and on a finished whole of \(1\).
The curves meet only when the parts are equal.
The mauve line is that even cut.}
\end{figure}

\subsection{Inversion of the pair}

Send \((x,2-x)\) to \(\bigl(1/x,\,1/(2-x)\bigr)\) for \(x\in(0,2)\).
The new means are
\[
\mathrm{AM}_{\mathrm{new}}=\frac{1}{x(2-x)},\qquad
\mathrm{GM}_{\mathrm{new}}=\frac{1}{\sqrt{x(2-x)}}.
\]

\begin{proposition}
\(\mathrm{AM}_{\mathrm{new}}=(\mathrm{GM}_{\mathrm{new}})^{2}\) everywhere on \((0,2)\), and \(\mathrm{AM}_{\mathrm{new}}\ge\mathrm{GM}_{\mathrm{new}}\) with equality only at \(x=1\).
\end{proposition}

\begin{proof}
The identity is the two formulae.
The inequality is Proposition~\ref{prop:amgm} on the inverted pair.
At \(x=1\) the pair is still \((1,1)\).
\end{proof}

Inversion around the unit appears here for the first time and will return whenever a close is written as a pair of inverse parts.
The lock \(x=1/x\) is the only place the two writings coincide; everywhere else the pair names each other and is not yet equal.

\begin{figure}[ht]
\centering
\includegraphics[width=0.92\textwidth]{drawing_03_inversion.png}
\caption{The same pair before and after inversion through the unit.
Equality remains only at \(x=1\).}
\end{figure}

\begin{lemma}[Midpoint]
For any segment \([A,B]\) the midpoint \(M=(A+B)/2\) is the unique point that is additively equidistant from the ends.
Every subsegment likewise admits a midpoint.
\end{lemma}

\begin{proof}
If \(M'\) is also additively equidistant from \(A\) and \(B\), then \(M'-A=B-M'\), hence \(M'=(A+B)/2=M\).
The same argument applies to any subsegment.
\end{proof}

In projective language the same cut is the harmonic conjugate of infinity:
\[
(A,B;M,\infty)=-1.
\]
On \([0,1]\) one has \(M=1/2\), and the harmonic permutations produce the family \(\{-1,1/2,2\}\).
Thus \(2=1/(1/2)\) already sits with inversion in the same set.
The mark \(2\) will later be the first prime, but that fact is not used here --- inversion already names it.

\begin{definition}[Primary ratio and prime ratio]
On the closed unit \([0,1]\), the unique primary, scale-invariant, \(\pm\)-equal self-description of location is the midpoint ratio \(1/2\).
The ``\(1\)'' in that writing is the closed unit itself.
The inverse of the primary ratio is \(2\).
Call \(2\) the \textbf{prime ratio}: it is the first mark inversion names from the primary cut, the first prime, and the only even prime.
Later primes are the same local irreducibility --- an arrival whose index does not factor through earlier finished counts.
The prime ratio is not a second cut of \([0,1]\).
It is the inverse of the only primary cut.
This is a naming choice, not a derivation.
That \(2\) is both the inverse of the primary cut and the first prime is a coincidence of small numbers worth marking, not a claim that primality follows from inversion.
No later result in this work depends on that convergence holding for any \(n\) other than \(2\).
\end{definition}

Those nested midpoints are the same primary ratio asked again: not new cuts, and not the prime ratio.

% ---------------------------------------------------------------------
\section{The circle of diameter 1}
% ---------------------------------------------------------------------

\subsection{When it closes}

Consider a circle whose diameter is exactly \(1\).
The radius is \(\tfrac12\).
The moment of closure is the integer \(1\) --- one finished turn --- and that integer is the same \(1\) as the opening question.

If a point runs on the rim at rate \(1\), then
\[
s=t,\qquad \theta=\frac{s}{1/2}=2t.
\]
One full turn is \(\theta=2\pi\), hence \(s=\pi\).
The finished turn is the integer \(1\); the length of the walk that made it is \(\pi\) diameters.
Those two statements are not in competition: one names the close, the other names the walk.

\begin{observation}[Live walk versus spelling]
Closure of a live walk is not an endpoint in the count of decimals.
It is arrival at HERE: a magnitude whose spelling has no last place.
\end{observation}

The tension is real, and it is not dissolved by saying ``\(\pi\) is the answer; stop asking which digit,'' nor is it an error to be corrected by preferring representation over geometry or geometry over representation.
Both respects are present: the walk closes, and the spelling does not.
They are two respects of one circumference.

\begin{figure}[ht]
\centering
\includegraphics[width=0.42\textwidth]{drawing_04_circle.png}
\caption{The paper circle. Diameter \(1\). Close at the integer \(1\).
Mauve is the even cut.}
\end{figure}

\subsection{Inverse parts of the close}

At that moment one may write \(\pi=1/\pi\) as a \emph{ratio of parts}, which is not a claim that the two numbers are equal.

\begin{proposition}
\label{prop:pi-parts}
\[
\pi\cdot\frac1\pi=1,\qquad
\mathrm{GM}\Bigl(\pi,\tfrac1\pi\Bigr)=1,\qquad
\mathrm{AM}\Bigl(\pi,\tfrac1\pi\Bigr)=\frac{\pi+1/\pi}{2}>1.
\]
\end{proposition}

\begin{proof}
The product is \(1\) by cancellation.
The geometric mean is \(\sqrt{1}=1\).
The arithmetic mean exceeds \(1\) by Proposition~\ref{prop:amgm}, since \(\pi\neq 1/\pi\).
\end{proof}

So \(1\) is the whole, and \(\pi\) and \(1/\pi\) are inverse parts of that whole.
On the inversion curve \(x\cdot(1/x)=1\), the equal case is the point \((1,1)\), and the parts of this circle are the point \(\bigl(\pi,\,1/\pi\bigr)\).
Equal, here, means related by inversion through the unit: the same writing will later be used for \(\Phi=1/\varphi\) and for \(1=1/1\), and nothing in the decimals has to match for the relationship to hold.

\begin{figure}[ht]
\centering
\includegraphics[width=0.92\textwidth]{drawing_05_parts.png}
\caption{The unit as whole, and the inverse parts \(\bigl(\pi,\,1/\pi\bigr)\) on \(xy=1\).}
\end{figure}

\subsection{Two halves that are not the same number}

Half the segment is \(\tfrac12\) and half the run is \(\pi/2\).
Both are honest halves of the same object, and they are not interchangeable quantities.

\begin{proposition}
\label{prop:two-halves}
On a circle of diameter \(d\),
\[
\frac{\text{half-walk}}{\text{half-diameter}}=\pi.
\]
In particular this is \((\pi/2)/(1/2)=\pi\) when \(d=1\), and \((n\pi/2)/(n/2)=\pi\) when \(d=n\).
\end{proposition}

\begin{proof}
The full walk is \(\pi d\).
Half the walk is \(\pi d/2\).
Half the diameter is \(d/2\).
The quotient is \(\pi\), independent of \(d\).
\end{proof}

This is the first rigidity: \(\pi\) is not recomputed at each size, because it is the locked ratio of the two halves of the template.

\begin{figure}[ht]
\centering
\includegraphics[width=0.42\textwidth]{drawing_06_two_halves.png}
\caption{Two halves of the paper circle: segment \(1/2\), walk \(\pi/2\).}
\end{figure}

% ---------------------------------------------------------------------
\section{The closed unit}
% ---------------------------------------------------------------------

The same object now needs a name that can travel with the count.

\begin{definition}[Closed unit]
The \textbf{closed unit} is one object in two writings.

As a unit it is \(1\): the lock \(x=1/x\), the geometric mean of every inversion pair \(\{x,1/x\}\), the place where the arithmetic mean falls onto that mean.
The primary ratio of that unit is \(1/2\).
The ``\(1\)'' in \(1/2\) is this unit.

As a closed package it is \(2\): the inverse of the primary ratio, named at count \(1\) and walked at count \(2\).
That length-\(2\) package is the first finished closed unit and the ruler used thereafter.
The prime ratio is this \(2\).

Arrival is comparison of readings at a finished step; neither mean knows arrival by itself.
\end{definition}

At count \(1\) the inverse of the midpoint of the finished unit \([0,1]\) names \(2\) without walking there.
The arithmetic mean of that pair is \(5/4\).
Until count \(2\) both ends of \([1/2,2]\) are not yet walked, so \(5/4\) is only a dated result.
At count \(2\) the pair is built and \(5/4\) becomes the ordinary midpoint of a walked segment.
The construction does not invent a location; it waits until both ends exist.

The closed unit is not only a list of marks. It is a rigid figure.
At finished \(n\), lift the own-origin to height \(n/2\).
The vertices
\[
(0,0),\qquad (n,0),\qquad \bigl(\tfrac n2,\tfrac n2\bigr)
\]
form a \(45^\circ\)--\(45^\circ\)--\(90^\circ\) triangle whose hypotenuse is the live segment itself.
The two slants have length \((n/2)\sqrt{2}\).
The upper semicircle of radius \(n/2\) centred at the own-origin passes through the apex; the half-turn of that circle is \(n\pi/2\).

Two lengths must not be confused.
The unit is \(1\) and its primary ratio is \(1/2\); the closed package is \(2\), with live length \(2\) and first half-turn \(\pi\).
At first closure the own-origin of the package is \(1\), which is already the lock; later packages have own-origin \(n/2\).
The cut of each package remains the primary ratio of that package, so the alphabet does not grow.

\begin{definition}[Master sequence]
The only allowed master is the successor count \(0,1,2,3,\ldots\).
One cannot go ahead in the count.
One finishes a step; then the next index exists; then one measures.
Privilege, if any, sits on the index of a completed step.
Geometry lives in the results.
Infinity is not a tick of the master sequence.
\end{definition}

After each finished \(n\) the measurements are forced:
\begin{align*}
\text{own-origin of }[0,n] &= n/2,\\
\text{inverse of the count} &= 1/n,\\
\text{cross-ratio }(0,1;n,1/n) &= -n.
\end{align*}

\begin{lemma}[Cross-ratio of the arrival package]
For every integer \(n\ge 2\),
\[
(0,1;n,1/n)=-n.
\]
\end{lemma}

\begin{proof}
\[
(0,1;n,1/n)
=\frac{(n-0)/(n-1)}{(1/n-0)/(1/n-1)}
=\frac{n/(n-1)}{(1/n)/((1-n)/n)}
=\frac{n}{n-1}\cdot(1-n)
=-n.
\]
\end{proof}

What the entrance has forced is small, and it is enough to walk: a finished unit, an even cut that is the only primary self-ratio, a circle that closes as an integer and walks as \(\pi\), and a count that is not allowed to go ahead.

% =====================================================================
\section{The integer count on the unit}
% =====================================================================

Each integer that follows is not a new object arriving from outside.
It is the same closed unit, asked again, at a larger scale --- the same ``where is \(1\)'' now posed as ``where is \(n\).''

Keep the whole as the unit \(1\), so that \(n\) means \(n\) parts of that \(1\).

\begin{proposition}
\label{prop:equal-split}
If \(n\) positive parts sum to \(1\), then \(\mathrm{AM}=1/n\), and \(\mathrm{GM}=\mathrm{AM}\) if and only if every part equals \(1/n\).
\end{proposition}

\begin{proof}
The arithmetic mean of parts summing to \(1\) is \(1/n\).
AM--GM gives \(\mathrm{GM}\le\mathrm{AM}\), with equality if and only if all parts are equal.
\end{proof}

Every integer also names a pair through the lock.

\begin{proposition}
\label{prop:inversion-pair}
For every positive \(n\),
\[
n\cdot\frac1n=1,\qquad
\mathrm{GM}\Bigl(n,\tfrac1n\Bigr)=1,\qquad
\mathrm{AM}\Bigl(n,\tfrac1n\Bigr)=\frac{n+1/n}{2}\ge 1,
\]
with equality in the last only at \(n=1\).
\end{proposition}

\begin{proof}
Immediate from cancellation and Proposition~\ref{prop:amgm}.
\end{proof}

The geometric mean of every such pair is the stopped mark \(1\), and the arithmetic mean meets that mark only at coincidence --- the same sentence as the first AM--GM drawing, now written along the count.

\begin{figure}[ht]
\centering
\includegraphics[width=0.92\textwidth]{drawing_08_two_walks.png}
\caption{Two walks inverted through \(1\), and the arithmetic mean sitting above a geometric mean that never leaves the lock.}
\end{figure}

\subsection{The local circle at integer \(n\)}

At finished count \(n\) the independent length sitting in the stock is the own-origin \(n/2\).
Place that length in the cross-section as the radius of a circle (equivalently: diameter equal to the arrived count \(n\)).
Name one further rigid move: a half-turn in that slice.
The resulting arc length is
\[
\frac{n\pi}{2}.
\]

With centre at the own-origin and radius \(n/2\), the upper semicircle passes through the apex of the \(45^\circ\)--\(45^\circ\)--\(90^\circ\) triangle of height \(n/2\).
The old lift height and the new local arc are two writings of one finished package.

Two measures of one segment must not be collapsed --- axis length \(n\) and half-walk \(n\pi/2\) --- and they lock at \(n=2\), where the length is \(2\) and the arc is \(\pi\).
Later circles are that first circle with a larger radius.
The interior spokes, the regular \(n\)-gon, the apothem \((n/2)\cos(\pi/n)\), the chord \(n\sin(\pi/n)\), and the central angle \(2\pi/n\) are all forced by the integer already on the table.

\begin{proposition}
\label{prop:half-walk}
At every integer \(n\ge 1\), half the local walk is \(n\pi/2\), and
\[
\frac{n\pi/2}{n/2}=\pi.
\]
\end{proposition}

\begin{proof}
The local diameter is \(n\).
Apply Proposition~\ref{prop:two-halves}.
\end{proof}

The ratio does not move.
Labels do.

\subsection{Two motions}

The frozen circle is a section.
The running count is a plan, and if height is taken to be walked arc it is an elevation; the conical helix of radius \(n/2\) is the oblique of that elevation.
Both writings are true, and neither is the other, which is why a single view cannot hold them and an architectural set can.

The temptation is to freeze the picture in order to have a place to stand.
The frozen picture is useful and accurate --- one-over-\(\pi\) is still \(0.318\ldots\) --- and the other writing is still available: \(\pi\) as parts against a whole of \(1\).
The cheap move is to say the circle is already closed.
The work is the count, and the walk is allowed.

\subsection{Rigidity}

Scale is a ratio of parts.
The picture does not change when \(n\) moves; labels and tick density do.
Run the count once and the measurement does not wander, which is why the construction can be measured at all: it is the same building at every scale.

That fact will matter later against a familiar objection, that infinity once admitted as a completed object changes the kind of the mark.
It does not.
Zoom is the same cut asked again.
If an architect puts the wrong scale on a drawing, the wall changes length and the building changes, but the drawing does not: the construction owns the relationship of parts, and scale is what you write in the title block.

\subsection{What every integer shares}

Collect the stock, dated by a finished \(n\ge 2\):
\begin{align*}
n\cdot(1/n)&=1,\\
\mathrm{GM}(n,1/n)&=1,\\
\mathrm{AM}(n,1/n)&=(n+1/n)/2,\\
\text{own-origin}&=n/2,\\
\text{inverse}&=1/n,\\
(0,1;n,1/n)&=-n,\\
\text{half-walk}&=n\pi/2,\\
\text{chord}&=n\sin(\pi/n),\\
\text{apothem}&=(n/2)\cos(\pi/n),\\
\text{central angle}&=2\pi/n.
\end{align*}

The alphabet does not grow.
The table through \(n=16\) is the stock at a depth large enough to see every letter; it belongs with the drawings at the end of this count, and it is not a search.

% ---------------------------------------------------------------------
\section{Vacuum and trace}
% ---------------------------------------------------------------------

Before the count is allowed to deepen, two respects of the same relation have to be named, or a name will be treated as an arrival.

\begin{definition}[Vacuum and trace]
\textbf{Vacuum:} the respect in which a relation may be named without being walked.
Maps and completions may appear at once.
The inversion \(x\mapsto 1/x\) is instant.
Infinity may be named.
Scaling toward \(0\) or \(\infty\) is free.
Vacuum here is not a physical vacuum and not an empty room.
It is the side of a relation that has no path-cost.

\textbf{Trace:} the same relation, if inscribed as positions, has a path.
With a finite speed of comparison there is a cost; the cost vanishes where a pair coincides and grows where vacuum scales freely.
\end{definition}

The same mathematical relation admits two respects, and naming is not the same as building.
One may write \(x=1/x\) in a single stroke; walking from \(x\) to \(1/x\) has a route and a duration.
The first is vacuum, the second is trace.
Collapse the two and a name is treated as an arrival --- privilege returns by the back door --- so neither respect is allowed to erase the other.

Vacuum may name a pair in one stroke.
Trace must walk both ends before a comparison is finished.
A finished comparison is a unit; an unfinished name is not.
That distinction is the whole of this section.

Rate is the other pole of the same pair, and the construction is scale-invariant even in rate.
Rate \(0\) is pause: one HERE, no next tick.
Any rate other than \(0\) is the count moving.
A day between integers, a year, a millisecond, the speed of light --- no tempo reaches a different kind of mark called \(\infty\), because infinity is still not reached as a trace.
The arrival package itself, \(n\) and \(1/n\) and the cross-ratio \(-n\), does not care what watch you used.

The walk around the paper circle is already one independent process: it has its own length, its own half-turn, and its own way of finishing a cycle.
Inversion is another, pairing every mark with a counterpart through the lock.
The even cut of a finished segment is a third, and it does not wait for either of the other two to exist.
Each of those processes carries its own ratios and can be followed without borrowing a scale from outside.
If each of them is casually called ``the unit,'' a name is treated as an arrival again.
They need a word that is not the unit.

\begin{definition}[Street and unit]
A \textbf{street} is an independent rigid process already owned by the closed unit: it carries its own ratios and its own infinity.
A \textbf{unit} is a finished relative comparison among streets.
Trace precedes measurement; measurement precedes unit.
The construction appoints neither a unit nor a vanishing.
\end{definition}

THERE is not a mark waiting on the ruler.
THERE is elected when a comparison arrives, and arrival makes a unit.
A unit the construction can seat sits on the \(1/2\) street --- the even cut that every finished segment already has.
Change the comparison and the unit changes; the street does not.

% ---------------------------------------------------------------------
\section{What stands}
% ---------------------------------------------------------------------

The entrance is finished enough to walk.
What it holds is not a list of refusals.
It is a small stock of things that have actually arrived, and a few things that have not been granted the job of arriving.

Forced, and available at every finished \(n\): the midpoint of a segment; inversion through the lock, with geometric mean \(1\) on every pair \(\{x,1/x\}\); the arithmetic mean sitting at or above that lock and meeting it only at coincidence; the inverse of \(\tfrac12\) equal to \(2\); two halves of the paper circle locking as \(\pi\); own-origin \(n/2\); inverse \(1/n\); the cross-ratio \((0,1;n,1/n)=-n\); the \(45^\circ\)--\(45^\circ\)--\(90^\circ\) triangle on every finished package.

Named, so they can be pointed at without being imported: closed unit, primary ratio, prime ratio, lock, paper circle, live walk against spelling, vacuum, trace, street, unit, master sequence.

Not granted: \(\pi\) as a constant carried in from outside the walk that produced it; a last decimal as closure of that walk; a completed run of the master sequence; a name treated as an arrival.

The frozen interior of the circle is useful, and it is accurate, and it is not allowed to cancel the live count.
A half-walk is the street measuring itself, not a later comparison wearing the costume of a finished unit.
Infinity is already interior to \((0,1)\); it is not a last index of the count, and no tempo makes it one.

What follows will only ask the same unit again --- at the next integer, through a second process on the same midline, and in the interiors those processes make when they work on each other.
The first sentence is still the sentence.
Where is \(1\)?
Point to it.

\end{document}
